% overlap-kp.tex — Chapter on the buffer-overlap Kotecky-Preiss
% certificate (rho_cap, K_cap) and the per-step Delta_cap budget.
%
% Linked code / Lean:
%   - computations/overlap_kp.py
%   - lean/YMGap/Sheaf/Gluing.lean (domGlue_preserves_KP)

\chapter{The Overlap KP Certificate and the
  \texorpdfstring{$\Dcap$}{Delta\_cap} Budget}
\label{ch:overlap-kp}

\noindent\emph{The buffer-overlap Kotecký--Preiss certificate
$(\rhocap, \Kcap)$ obtained by gluing the small-field and
large-field certificates via the domination construction of the
sheaf framework, and the per-step buffer growth $\Dcap(b, g_k) =
\Kcap(k+1) - \Kcap(k)$ that feeds the $\sigh$-budget.}

This chapter combines the small-field certificate of
\Cref{ch:small-field-kp} with the large-field certificate of
\Cref{ch:large-field-kp} via the domination gluing of
\Cref{ch:sheaf-framework} to produce the buffer-overlap certificate
$c_k^{\cap}$.  The constants relationship
$\rhocap = \min(\rhosf, \rholf)$ and $\Kcap \le \Ksf + \Klf$ is the
framework-level claim, and the per-step buffer growth $\Dcap$ — the
quantity that feeds the $\sigh$-budget of
\Cref{ch:sigma-hat-trajectory} — is defined as the discrete
difference $\Kcap(k+1) - \Kcap(k)$.  We also state the empirical
looseness bound that motivates the chaining tightening of
\Cref{ch:chaining-estimate}.

\section{The overlap certificate}
\label{sec:overlap-certificate}

\marginkey{Domination gluing}%

Recall from \Cref{ch:sheaf-framework},
Proposition~\ref{prop:summability-preservation}, the
\emph{domination-gluing} construction: given KP certificates
$(\rhosf, \Ksf)$ on $\Omksf$ and $(\rholf, \Klf)$ on $\Omklf$, the
buffer overlap region $\Omkcap$ inherits a glued certificate.

\begin{theorem}[Domination gluing on the buffer]
\label{thm:overlap-glue}
Suppose at scale $k$ both certificates converge: $\rhosf(g_k) >
\mu_{\mathrm{conn}}$ (\Cref{ch:small-field-kp},
Theorem~\ref{thm:ksf-convergence}) and
$\muplaq\,Z_t^{\mathrm{upper}}(g_k) < 1$ (\Cref{ch:large-field-kp},
Theorem~\ref{thm:klf-convergence}).  Then the buffer overlap region
$\Omkcap$ carries a Kotecký--Preiss certificate $(\rhocap, \Kcap)$
with constants
\begin{align}
  \label{eq:rho-cap}
  \rhocap(g_k) &\;:=\; \min\bigl(\rhosf(g_k),\, \rholf(g_k)\bigr), \\
  \label{eq:k-cap}
  \Kcap(g_k) &\;\le\; \Ksf(g_k) + \Klf(g_k) .
\end{align}
\end{theorem}

\begin{proof}[Proof sketch]
On the buffer overlap, every polymer activity is dominated by both
the small-field and the large-field activity bound; therefore its
maximum is dominated by either, and the slower of the two decay
rates controls.  The summability constant inherits the union of
both contributions, giving $\Kcap \le \Ksf + \Klf$ by the triangle
inequality applied to the polymer-cluster expansion on the
disjoint union $\Omksf \sqcup \Omklf \supseteq \Omkcap$.  The
detailed argument is in \Cref{ch:sheaf-framework}; the present
chapter applies it at the $\SUtwo$ regime-specific level.
\end{proof}

\begin{remark}[Saturation of the union bound]
\label{rem:K-cap-saturation}
The bound $\Kcap \le \Ksf + \Klf$ is the worst-case union bound; it
saturates when the small-field and large-field activity bounds are
\emph{disjoint} in polymer support.  In practice they overlap in
the buffer annulus, so the actual $\Kcap$ is strictly less than
$\Ksf + \Klf$ — the looseness factor of
\Cref{ch:overlap-measure}, Definition~\ref{def:looseness}, recurs
here, with the same chaining-tightening remedy of
\Cref{ch:chaining-estimate}.
\end{remark}

\begin{definition}[Overlap KP certificate]
\label{def:overlap-certificate}
The \emph{buffer-overlap Kotecký--Preiss certificate at scale $k$}
is the triple
\begin{equation}
  \label{eq:cap-certificate}
  c_k^{\cap} \;:=\;
  \bigl(\rhocap(g_k),\, \Kcap(g_k),\, \mathrm{sat}_{\cap}\bigr)
  \;\in\; \mathrm{KPCertificate}(V, P_{\cap}),
\end{equation}
where $\mathrm{sat}_{\cap}$ is the proof of
\[
  P_{\cap}(\rhocap, \Kcap) :\Leftrightarrow
  \rhocap > 0 \;\wedge\; \Kcap < \infty
\]
delivered by \cref{thm:overlap-glue}.
\end{definition}

\section{The per-step buffer growth
  \texorpdfstring{$\Dcap$}{Delta\_cap}}
\label{sec:delta-cap}

\marginkey{Buffer-$K$ growth per RG step}%

The buffer-overlap summability constant $\Kcap(g_k)$ is
\emph{coupling-dependent}: as $g_k$ flows along the RG trajectory,
$\Kcap$ changes.  The per-step growth of $\Kcap$ is the
buffer-budget quantity $\Dcap$ used by the $\sigh$-transport of
\Cref{ch:sigma-hat-trajectory}.

\begin{definition}[Per-step buffer growth]
\label{def:delta-cap}
Let $b > 1$ be the block-spin factor of \Cref{ch:block-spin-morphism}.
The \emph{per-step buffer growth at scale $k$} is
\begin{equation}
  \label{eq:delta-cap}
  \Dcap(b, g_k) \;:=\; \max\bigl(0,\;\; \Kcap(g_{k+1}) - \Kcap(g_k)\bigr),
\end{equation}
where $g_{k+1}$ is the next-scale coupling determined by the
one-loop $\beta$-function (\Cref{ch:uv-trajectory}).
\end{definition}

The $\max(0, \cdot)$ truncation reflects the buffer growth being
monotone non-decreasing across the block-spin step: in the UV-direction
$g$ flows downward, both $\Ksf$ and $\Klf$ flow downward,
and so does $\Kcap$, giving $\Dcap = 0$.  In the IR-direction
flow $g$ grows and $\Kcap$ grows; $\Dcap > 0$ captures the
incremental buffer cost.

When either endpoint's $\Kcap$ is non-finite (the
non-convergent regime), $\Dcap = \infty$ by convention, propagating
the non-convergence honestly through the downstream budget.

\begin{mainfigure}
\centering
\includegraphics[width=\textwidth]{delta-cap-budget}
\caption{Per-step buffer growth $\Dcap(b, g_k)$ across the canonical
$\SUtwo$ RG trajectory at $b = 2$.  The geometric decay matches
the analytic UV-direction rate $r_{\mathrm{uv}}^\infty \approx
0.976$ used by the global $\sigh$-budget closure of
\Cref{ch:sigma-hat-trajectory}.}
\label{fig:delta-cap-budget}
\end{mainfigure}

\section{Empirical looseness and the gaussian-sharpness
  conjecture}
\label{sec:k-cap-looseness}

\Cref{rem:K-cap-saturation} flagged the union bound $\Kcap \le
\Ksf + \Klf$ as worst-case.  Empirically, on the $\SUtwo$ canonical
convergent trajectory we observe the union-bound looseness factor
\begin{equation}
  \label{eq:k-looseness}
  L_K(g_k) \;:=\; \frac{\Ksf(g_k) + \Klf(g_k)}{\max\bigl(\Ksf(g_k),\, \Klf(g_k)\bigr)}
  \;=\;
  1 + \frac{\min(\Ksf, \Klf)}{\max(\Ksf, \Klf)}
\end{equation}
stays comfortably below $1.2$ across the convergent regime, with
peak $\approx 1.18$ at $g = 0.35$.

\begin{conjecture}[Gaussian sharpness; \Cref{ch:overlap-measure},
  Conjecture~\ref{conj:gaussian-sharpness}, restated]
\label{conj:K-cap-sharpness}
There exists $C \le 2$ such that on every trajectory in the
joint-convergent regime, $L_K(g_k) \le C$.
\end{conjecture}

The conjecture is the gaussian-sharpness analog of
\Cref{conj:gaussian-sharpness}, lifted to the $\Kcap$-summability
level.  Empirically the bound is much tighter than $2$, but the
universal constant $C \le 2$ remains the manuscript's open
theoretical target.

\begin{infobox}[main][title=Why $\Ksf$ dominates]
The small-field and large-field certificates approach divergence
from opposite directions: $\Ksf \to \infty$ as $g \to
g_{\max}^{\mathrm{sf}}$ from below, while $\Klf \to 0$ as $g \to 0$
(the activity bound trivializes).  At the joint-convergent regime
boundary, $\Ksf \gg \Klf$, and the union-bound looseness factor
collapses to $1 + \Klf/\Ksf \approx 1 + \mathrm{small}$.
\end{infobox}

\begin{mainfigure}
\centering
\includegraphics[width=\textwidth]{conjecture-7.4-looseness}
\caption{The empirical union-bound looseness factor
$L_K(g_k) = 1 + \min(\Ksf, \Klf)/\max(\Ksf, \Klf)$ of
\cref{eq:k-looseness} across the canonical $\SUtwo$ convergent
regime.  $L_K$ peaks at $\approx 1.18$ near $g = 0.35$, far below
the gaussian-sharpness conjecture bound $C \le 2$.  At the
canonical pilot $g_0 = 0.30$, $L_K \approx 1.01$ — the union bound
is essentially saturated.}
\label{fig:conjecture-7.4-looseness}
\end{mainfigure}

\section{Convention and the $\Dcap$ budget in
  \texorpdfstring{$\Qp$}{Q'}}
\label{sec:overlap-Q-rational}

Both $\Ksf$ and $\Klf$ are constructively $\Qp$-rational at every
pilot operating point (\Cref{ch:small-field-kp,ch:large-field-kp},
respectively), so $\Kcap \le \Ksf + \Klf$ is $\Qp$-rational by
closure under addition.  $\Dcap(b, g_k) = \Kcap(g_{k+1}) -
\Kcap(g_k)$ is therefore also $\Qp$-rational, provided the
trajectory's $\{g_k\}$ values are $\Qp$-rational — which they are at
every finite-$K$ pilot truncation of \Cref{ch:sigma-hat-trajectory}.

\section{Use throughout the manuscript}
\label{sec:overlap-kp-downstream}

\Cref{ch:sigma-hat-trajectory} feeds $\Dcap(b, g_k)$ into the
$\sigh$-transport derivation: the per-step degradation $\eta_k$ is
proportional to $\Dcap$, with proportionality constant
$\cpolymer / \Amin = 2/4 = 1/2$.  The full chain
$\Ksf, \Klf \to \Kcap \to \Dcap \to \eta_k \to \sum_k \eta_k <
\sigma_0$ closes the budget at the $\Qp$-rational level at finite
truncation $K$, supplying the constructive mass-gap lower bound
$\massgap$ of \Cref{ch:mass-gap-theorem}.

\section{Code provenance}
\label{sec:overlap-kp-codeanchor}

\begin{itemize}
  \item Python implementation:
    \codepath{computations/overlap\_kp.py}.  Key entry points:
    \codepath{compute\_overlap\_kp(g\_k, t\_choice, scale\_k,
    spectral\_gap\_coeff, mu\_plaq)} returns an
    \codepath{OverlapKP} record carrying $\rhocap$, $\Kcap$, both
    underlying regime certificates, and a provenance dictionary;
    \codepath{delta\_cap(g\_k, g\_k\_next, t\_choice)} computes the
    per-step buffer growth of \cref{def:delta-cap};
    \codepath{union\_bound\_looseness\_estimate(g\_k, t\_choice)}
    surfaces the empirical $L_K$ of \eqref{eq:k-looseness}.
    Tests in \codepath{tests/test\_overlap\_kp.py}.
  \item Lean formalization:
    \codepath{lean/YMGap/Sheaf/Gluing.lean::domGlue\_preserves\_KP}
    is the type-level statement of \cref{thm:overlap-glue}.  The
    Lean theorem carries part~(a) of the open-obligation ledger
    entry $[P2.1]$; part~(b) is \cref{conj:K-cap-sharpness},
    currently open.
  \item Visualizations:
    \codepath{visualizations/pilot/su2/delta-cap-budget.svg}
    plots $\Dcap(b, g_k)$ across the canonical trajectory;
    \codepath{visualizations/pilot/su2/conjecture-7.4-looseness.svg}
    plots $L_K(g_k)$.
\end{itemize}
